\section{Rank-distance experiments and switching constraints}
\label{sec:path-rank-distance-structure}

The following observations concern the unresolved quantitative target
\[
 P_{uv}(N)\le2^{b-1}+\operatorname{poly}(b)\,2^{(b+\ell)/2},
 \qquad \ell=\operatorname{dist}_N(u,v).
\]
The Horak--Kim construction in Proposition~\ref{prop:hk-short-path-obstruction} rules out replacing the polynomial by an absolute constant.
The observations below do not prove this target.

\begin{lemma}[An exact path-complement identity]
\label{lem:path-complement-partition}
Let $N$ be a normalized two-pole of rank $b$, with terminals $u,v$.
Then
\[
 \sum_{Q\text{ a simple }u\text{--}v\text{ path}}
 2^{\beta(N-V(Q))}=2^b.
\]
In particular,
\[
 P_{uv}(N)-2^{b-1}
 =\sum_Q\left(1-2^{\beta(N-V(Q))-1}\right).
\]
\end{lemma}

\begin{proof}
Every $\{u,v\}$-join has selected degree one at the two terminals and selected degree zero or two elsewhere.
Thus it consists of one terminal path $Q$ and an even subgraph of $N-V(Q)$.
Conversely, every such pair gives a unique terminal join.
There are $2^b$ terminal joins, and the even space of the complementary graph has $2^{\beta(N-V(Q))}$ states.
This proves the first identity.
Subtract half that identity from the number of paths to obtain the second.
\end{proof}

Only paths with forest complement make a positive contribution in the second identity.
Discarding the negative contributions is potentially costly: the positive class includes Hamilton terminal paths.
A proof at the half exponent must retain sufficient cancellation or use another structural restriction.

\begin{lemma}[A switching upper bound for the excess]
\label{lem:path-switching-excess}
Let $\mathcal A$ be the set of terminal paths, viewed as a subset of the affine terminal-join space of size $2^b$.
For an even edge set $F$, write
\[
 R_F=|\mathcal A\cap(\mathcal A\mathbin\triangle F)|.
\]
Then
\[
 P_{uv}(N)-2^{b-1}\le R_F/2.
\]
Consequently, if the path density exceeds one half, every even edge set is the symmetric difference of two terminal paths.
More quantitatively, it has at least $2P_{uv}(N)-2^b$ ordered representations of that form.
\end{lemma}

\begin{proof}
Both $\mathcal A$ and its translate lie in the same affine space and have the same size.
Inclusion--exclusion gives $R_F\ge2|\mathcal A|-2^b$.
Each member of the intersection determines its unique translated partner.
\end{proof}

Thus the rank-distance target would follow by finding one even set $F$ with
\[
 R_F\le\operatorname{poly}(b)\,2^{(b+\ell)/2}.
\]
This is a sufficient condition, not an established estimate.
An arbitrary subset of a binary affine space need not have such a sparse overlap with any translate; graph structure is essential.

\begin{proposition}[A necessary threading condition for a spanning cycle packing]
\label{prop:path-factor-threading}
Let $N$ be a normalized two-pole, let $G=N+uv$, and write $e$ for the added edge.
Suppose $F\subseteq E(N)$ is a spanning union of vertex-disjoint cycles.
Contract each cycle of $F$ in $G$ and delete the contracted cycle edges, retaining loops and parallel edges.
If $P_{uv}(N)>2^{b-1}$, the resulting quotient has a connected spanning even subgraph containing the image of $e$.
\end{proposition}

\begin{proof}
The switching lemma supplies terminal paths $Q_1,Q_2$ with $Q_1\mathbin\triangle Q_2=F$.
Put $C_i=Q_i\cup\{e\}$ and $M=C_1\cap C_2$.
Since $F$ spans every vertex and has degree two at every vertex, $E(G)\setminus F$ is a perfect matching.
Hence $M$ is a subset of that matching and contains $e$.

Every cycle of $F$ contains an edge of each $C_i$.
Otherwise it is a cycle contained entirely in one $C_i$, forcing that $C_i$ to equal it, contrary to $e\in C_i$.
The image of the connected cycle $C_1$ after contraction therefore meets every quotient vertex and consists exactly of the edges of $M$.
It is connected and spanning.
At each contracted cycle, a cycle crosses its boundary an even number of times.
Thus the image of $M$ is even and contains $e$.
\end{proof}

The quotient condition can rule out path density above one half for a specified spanning cycle packing.
It does not yet supply a bound on the positive excess when the condition holds.

\subsection{Finite catalogue checks}

The program \texttt{check\_path\_rank\_deficit.cpp} enumerates the cycle space of every connected simple cubic graph in the catalogue through order twenty.
For each edge $e=uv$, cycles containing $e$ count terminal paths in $G-e$ exactly.
A separate breadth-first search gives the terminal distance.
The catalogue contains 556,471 graphs in this range.
The largest values at each terminal rank $b$ of
\[
 \left(P/2^b-\frac12\right)2^{(b-\ell)/2}
\]
are listed below.
\[
\begin{array}{c|rrrrrrrrr}
 b&2&3&4&5&6&7&8&9&10\\ \hline
 \text{maximum}&.500000&.530330&.625000&.687500&.839689&1.000000&1.259534&1.421875&1.800913
\end{array}
\]
At $b=10$, the maximum has $\ell=3$ and $P=675$.
The direct path-search program \texttt{check\_path\_rank\_deficit\_witnesses.py} independently verifies all 61 distinct recorded extremal witnesses, using no cycle-space enumeration or third-party package.
Its result is \texttt{path-rank-deficit-witness-checks.json}.
The full catalogue results are \texttt{path-rank-deficit-catalog-results.json}.
These checks are finite evidence only.

A tempting stronger decomposition would assert that every above-half path family has a full slice obtained by fixing one edge coordinate.
The edge-deleted Petersen graph is a finite obstruction to that assertion.
Its rank is five, its terminal distance is four, and it has 28 paths among 32 joins.
For each of the fourteen remaining edges, both coordinate values occur among its four bad joins.
Thus none of those slices is full.
This does not refute the rank-distance target, since $b-\ell=1$ in this example.
